\chapter{Especificación de requerimientos}

\section{Alcances y limitaciones}

\subsection{Alcances}
El presente proyecto abarca el diseño, desarrollo, verificación y empaquetado del módulo móvil de la plataforma \textbf{QuimioCare}, enfocado exclusivamente en el acompañamiento ambulatorio de pacientes adultos en tratamiento de \textbf{quimioterapia}. El alcance específico comprende:

\begin{enumerate}
    \item \textbf{Aplicación móvil Android nativa:} Desarrollo en React Native con Expo SDK 57 y TypeScript, generando una \textit{Development Build} y compilación \textit{Release} demostrativa independiente con identificador de paquete \texttt{cl.ubiobio.quimiocare}.
    \item \textbf{Arquitectura Offline-First estricta:} Persistencia relacional en el dispositivo mediante SQLite (\texttt{expo-sqlite}) con migraciones versionadas (\texttt{001} a \texttt{008}), transacciones ACID y modo WAL, garantizando la captura y consulta ininterrumpida de signos vitales, síntomas libres, cuestionarios médicos asignados y contexto clínico sin requerir conexión a internet.
    \item \textbf{Aislamiento de perfiles de prueba:} Soporte determinista para dos pacientes sintéticos de demostración (Paciente A y Paciente B), garantizando que las bases de datos, copias de archivos locales y recordatorios mantengan estricta separación de pertenencia.
    \item \textbf{Gestión de copias educativas locales:} Descarga, almacenamiento atómico y lectura sin conexión de guías e infografías educativas en almacenamiento privado (\texttt{expo-file-system}), con reconciliación bidireccional de estados y visualización en línea de contenidos audiovisuales alojados en YouTube.
    \item \textbf{Motor nativo de reportes clínicos:} Consolidación histórica de mediciones y síntomas en ventanas de 7, 30 y 90 días, con generación y exportación local de documentos PDF (mediante \texttt{expo-print} con plantillas sanitizadas) y planillas de cálculo Excel (\texttt{.xlsx} con SheetJS), incorporando el código único de vinculación del paciente para su presentación en controles médicos.
    \item \textbf{Coordinador local de recordatorios:} Planificación nativa de notificaciones de medicamentos y controles mediante \texttt{expo-notifications} en una ventana móvil de 14 días (hasta 128 trabajos lógicos), operativa de forma autónoma sin la aplicación en ejecución.
    \item \textbf{Frontera de integración diferida:} Registro atómico de cada evento clínico en una cola local de envío (\texttt{cola\_envio}) con UUID idempotente, preparada para su futura transmisión REST hacia la pasarela o el simulador central.
\end{enumerate}

\subsection{Limitaciones}

\subsubsection{Limitaciones de Hardware y Software}
\begin{enumerate}
    \item \textbf{Plataforma de destino:} El desarrollo y la validación nativa se acotan a dispositivos móviles con sistema operativo Android versión 12 (API 31) o superior. El ecosistema iOS queda fuera del alcance de este ciclo.
    \item \textbf{Dependencia de conectividad para video:} Los recursos audiovisuales alojados en YouTube requieren conexión a internet para su reproducción externa; la aplicación no descarga ni almacena localmente archivos de video para respetar derechos de autor y políticas de almacenamiento.
    \item \textbf{Procesamiento de lenguaje natural y modelos generativos:} La aplicación no ejecuta modelos de lenguaje (LLM) ni motores RAG en el hardware del dispositivo. El asistente opera en el cliente mediante un controlador con estados deterministas y disclaimers visibles; la generación avanzada de respuestas queda delegada al servicio central.
\end{enumerate}

\subsubsection{Limitaciones Operativas y Clínicas}
\begin{enumerate}
    \item \textbf{Exclusión de radioterapia:} El módulo excluye flujos, tablas y contenidos de radioterapia, debido a que el Hospital Clínico Herminda Martín no dispone de equipamiento de radioterapia en Chillán (derivando pacientes a Concepción) y a que su dinámica asistencial difiere radicalmente de las toxicidades agudas de la quimioterapia sistémica.
    \item \textbf{Principio de no intervención clínica autónoma:} La aplicación es una herramienta de registro, educación y acompañamiento; no emite diagnósticos, no sugiere cambios de dosis ni reemplaza la evaluación médica presencial de urgencias.
    \item \textbf{Autoridad en la definición de cuestionarios:} Los formularios clínicos (tales como QLQ-C30 y ESAS) son diseñados y asignados exclusivamente por el equipo de salud en el backend central; el paciente responde las preguntas vigentes pero no edita plantillas clínicas.
\end{enumerate}

\section{Planificación de reuniones con el cliente y levantamiento}
El levantamiento de requerimientos se fundamentó en sesiones técnicas y entrevistas cualitativas con el personal clínico de la Unidad de Quimioterapia Ambulatoria del Hospital Clínico Herminda Martín (HCHM) de Chillán, junto con la revisión del Plan Nacional de Cáncer y la coordinación con el equipo del módulo web y de inteligencia clínica. Dichas interacciones permitieron identificar las siguientes directrices esenciales:
\begin{itemize}
    \item Priorizar la simpleza de uso y minimizar la sobrecarga cognitiva mediante la heurística de \textit{máximo 3 interacciones} para alcanzar cualquier función principal.
    \item Asegurar funcionamiento completo sin señal móvil para dar cobertura a la alta ruralidad de la Región de Ñuble.
    \item Incorporar de forma obligatoria avisos de responsabilidad clínica y la exportación de reportes físicos para el control médico presencial.
\end{itemize}

\section{Requerimientos funcionales}

A continuación se presentan los requerimientos funcionales del módulo móvil agrupados por subsistema:

% --- TABLA REGISTRO CLÍNICO Y VITALES ---
\begin{table}[H]
\centering
\renewcommand{\arraystretch}{1.3}
\setlength{\tabcolsep}{6pt}
\begin{tabularx}{\textwidth}{|c|p{3.5cm}|X|c|}
\hline
\textbf{ID} & \textbf{Nombre} & \textbf{Descripción} & \textbf{Prioridad} \\ \hline
RF\_01 & Registrar signos vitales & El sistema debe permitir registrar mediciones de presión arterial (sistólica/diastólica), pulso cardíaco, temperatura axilar y saturación de oxígeno, admitiendo entradas parciales y guardándolas de forma atómica en SQLite. & Alta \\ \hline
RF\_02 & Registrar síntomas diarios & El sistema debe permitir redactar una descripción libre de síntomas y malestares de hasta 2.000 caracteres, registrando la fecha y hora exacta sin requerir conexión a internet. & Alta \\ \hline
RF\_03 & Responder cuestionarios médicos & El sistema debe presentar formularios y encuestas médicas asignadas (ej. QLQ-C30, ESAS), permitiendo guardar borradores locales editables y confirmar entregas inmutables. & Alta \\ \hline
RF\_04 & Consultar historial clínico & El sistema debe mostrar el historial cronológico de mediciones y síntomas con filtros temporales de 7, 30 y 90 días y cálculo de estadísticas descriptivas. & Alta \\ \hline
\end{tabularx}
\caption{Requerimientos funcionales – Registro clínico y seguimiento}
\label{tab:rf_clinico}
\end{table}

% --- TABLA EDUCACIÓN Y ASISTENCIA ---
\begin{table}[H]
\centering
\renewcommand{\arraystretch}{1.3}
\setlength{\tabcolsep}{6pt}
\begin{tabularx}{\textwidth}{|c|p{3.5cm}|X|c|}
\hline
\textbf{ID} & \textbf{Nombre} & \textbf{Descripción} & \textbf{Prioridad} \\ \hline
RF\_05 & Consultar catálogo educativo & El sistema debe organizar recursos educativos en categorías temáticas sobre quimioterapia, cuidados y nutrición. & Alta \\ \hline
RF\_06 & Descarga y lectura de copias offline & El sistema debe permitir guardar copias locales de infografías y guías en almacenamiento persistente para lectura sin red, gestionando su eliminación y espacio utilizado. & Alta \\ \hline
RF\_07 & Visualizar videos de YouTube & El sistema debe permitir abrir enlaces directos a videos informativos en la aplicación de YouTube o navegador externo cuando exista conexión a internet. & Media \\ \hline
RF\_08 & Interacción con asistente conversacional & El sistema debe ofrecer una interfaz de chat asistivo con estados deterministas, cancelación de solicitudes y aviso permanente de responsabilidad médica. & Media \\ \hline
\end{tabularx}
\caption{Requerimientos funcionales – Educación y Asistencia}
\label{tab:rf_educacion}
\end{table}

% --- TABLA REPORTES Y RECORDATORIOS ---
\begin{table}[H]
\centering
\renewcommand{\arraystretch}{1.3}
\setlength{\tabcolsep}{6pt}
\begin{tabularx}{\textwidth}{|c|p{3.5cm}|X|c|}
\hline
\textbf{ID} & \textbf{Nombre} & \textbf{Descripción} & \textbf{Prioridad} \\ \hline
RF\_09 & Generar reporte PDF local & El sistema debe construir un documento PDF imprimible con las mediciones y síntomas del período seleccionado, incluyendo identificación y código de vinculación. & Alta \\ \hline
RF\_10 & Exportar reporte a Excel & El sistema debe generar una planilla de cálculo (.xlsx) estructurada con los registros clínicos del paciente para su procesamiento o auditoría médica. & Alta \\ \hline
RF\_11 & Compartir reportes externamente & El sistema debe invocar la hoja de compartición nativa de Android para enviar o guardar los archivos PDF y Excel generados. & Alta \\ \hline
RF\_12 & Gestionar recordatorios locales & El sistema debe permitir configurar horarios de medicación y controles, planificando notificaciones nativas en una ventana móvil de 14 días. & Alta \\ \hline
RF\_13 & Bandeja de avisos del paciente & El sistema debe proveer una vista centralizada de los avisos programados y pasados con acceso directo a la acción asociada. & Media \\ \hline
\end{tabularx}
\caption{Requerimientos funcionales – Reportes y Recordatorios}
\label{tab:rf_reportes_recordatorios}
\end{table}

% --- TABLA IDENTIDAD Y CONFIGURACIÓN ---
\begin{table}[H]
\centering
\renewcommand{\arraystretch}{1.3}
\setlength{\tabcolsep}{6pt}
\begin{tabularx}{\textwidth}{|c|p{3.5cm}|X|c|}
\hline
\textbf{ID} & \textbf{Nombre} & \textbf{Descripción} & \textbf{Prioridad} \\ \hline
RF\_14 & Aislamiento de paciente (A/B) & El sistema debe permitir alternar entre pacientes sintéticos de demostración asegurando la separación estricta de sus registros y configuraciones. & Alta \\ \hline
RF\_15 & Visualizar código de vinculación & El sistema debe mostrar en pantalla el código alfanumérico único para vincular el dispositivo con el equipo médico en controles presenciales. & Alta \\ \hline
RF\_16 & Configurar apariencia accesible & El sistema debe permitir seleccionar entre cuatro paletas de color de alto contraste (Turquesa, Azul, Violeta, Ámbar) y modo claro/oscuro. & Media \\ \hline
RF\_17 & Registrar consentimiento versionado & El sistema debe presentar la versión vigente de términos y registrar el otorgamiento o revocación como eventos con trazabilidad local. & Alta \\ \hline
\end{tabularx}
\caption{Requerimientos funcionales – Identidad y Configuración}
\label{tab:rf_identidad}
\end{table}

\section{Requerimientos no funcionales}

Los requerimientos no funcionales definen los criterios de calidad, arquitectura, rendimiento y seguridad que rigen el sistema:

\begin{table}[H]
\centering
\renewcommand{\arraystretch}{1.3}
\setlength{\tabcolsep}{6pt}
\begin{tabularx}{\textwidth}{|c|p{3.2cm}|X|c|}
\hline
\textbf{ID} & \textbf{Criterio} & \textbf{Descripción y Métrica de Aceptación} & \textbf{Prioridad} \\ \hline
RNF\_01 & Usabilidad y Carga Cognitiva & Interfaz optimizada con controles de $\ge 48$\,dp, tipografía legible con soporte para texto ampliado del sistema y máximo de 3 interacciones de navegación para alcanzar cualquier función principal. & Crítica \\ \hline
RNF\_02 & Disponibilidad Offline-First & Todas las funciones de captura (vitales, síntomas, encuestas), lectura de historial y consulta de copias guardadas deben operar al 100\% sin conexión a internet. & Crítica \\ \hline
RNF\_03 & Integridad y Persistencia ACID & Operaciones relacionales en SQLite bajo modo WAL y claves foráneas, asegurando atomicidad entre el guardado del dato clínico y la inserción del evento en la cola de salida. & Crítica \\ \hline
RNF\_04 & Aislamiento Estricto de Datos & Prohibición absoluta de filtración o herencia de información entre perfiles; el Paciente A no puede leer, modificar ni eliminar datos pertenecientes al Paciente B. & Crítica \\ \hline
RNF\_05 & Seguridad y Sanitización & Confinamiento en sandbox privado del dispositivo, sanitización estricta contra inyección de HTML en plantillas PDF y protección de datos sensibles. & Alta \\ \hline
RNF\_06 & Rendimiento de Exportación & La generación y persistencia física de reportes PDF y Excel para hasta 90 días de historial debe completarse en menos de 3.000\,ms en hardware Android de referencia. & Alta \\ \hline
RNF\_07 & Compatibilidad de Plataforma & Compilación como aplicación Android nativa compatible con Android 12 (API 31) en adelante, gestionando permisos en tiempo de ejecución. & Alta \\ \hline
RNF\_08 & Responsabilidad Clínica & Exhibición permanente e ineludible de avisos que aclaren la naturaleza de acompañamiento de la app, instruyendo acudir a urgencias ante signos críticos. & Crítica \\ \hline
\end{tabularx}
\caption{Requerimientos no funcionales del sistema}
\label{tab:rnf}
\end{table}

\section{Matriz de trazabilidad}

La siguiente matriz vincula los Objetivos Específicos (OE) definidos en el Capítulo 3 con los requerimientos funcionales implementados:

\begin{table}[H]
\centering
\renewcommand{\arraystretch}{1.3}
\setlength{\tabcolsep}{6pt}
\begin{tabularx}{\textwidth}{|c|c|X|l|}
\hline
\textbf{Objetivo} & \textbf{Requerimientos} & \textbf{Módulos Asociados} & \textbf{Ámbito de Solución} \\ \hline
OE1 & RF\_01, RF\_02, RF\_05, RF\_08, RF\_16 & Interfaz, Vitales, Educación, Accesibilidad & Presentación accesible y captura en $\le 3$ toques \\ \hline
OE2 & RF\_01, RF\_02, RF\_08 & Registro Clínico, Asistencia & Monitoreo y advertencia ante síntomas críticos \\ \hline
OE3 & RF\_01, RF\_02, RF\_03, RF\_04, RF\_06, RF\_14 & SQLite Local, Repositorios, Sandboxing & Arquitectura offline-first y persistencia relacional \\ \hline
OE4 & RF\_08, RF\_12, RF\_13 & Gateway, Notificaciones, Sesiones & Coordinación de avisos locales y pasarela \\ \hline
OE5 & RF\_09, RF\_10, RF\_11, RF\_15 & Motor de Reportes (PDF/Excel), Vinculación & Exportación física de datos y vinculación médica \\ \hline
OE6 & RNF\_01 a RNF\_08 & Infraestructura de Pruebas, Vitest, CI & Verificación de robustez, aislamiento y rendimiento \\ \hline
\end{tabularx}
\caption{Matriz de trazabilidad – Objetivos Específicos vs. Requerimientos}
\label{tab:trazabilidad}
\end{table}